Ansible instead uses scripts called \texttt{Playbooks} to describe the target state of the machine, which are usually kept idempotent to ensure that re-running the same \texttt{Playbook} does not alter the state of the target machine if it already is in the requested state.
This section is going to give an introduction to Ansible and its components to give a good understanding of how Ansible works and what its requirements for execution are.
The section will only focus on components that are going to be relevant to the implementation of the operator.
\parencite{ansibledocsindex}

\subsection{The Ansible inventory}
The Ansible inventory is what tells Ansible what hosts to target, how they are grouped and what variables to use for them when evaluating \texttt{Playbooks} or \texttt{Roles}.
It is written in a file or a predetermined directory structure.
This thesis will build Ansible inventories using the directory layout, allowing easier manageability of scaling environments, as per \cite{ansibleinventoryintro}.

An Ansible inventory is separated into \texttt{Groups} and \texttt{Hosts}, both of which can have variables.
\texttt{Groups} are a way to group multiple \texttt{Hosts} into a logical unit, which can be targeted by a \texttt{Playbook}.
A \texttt{Group} is declared simply by creating a key in the YAML file that represents the inventory at root level.
\texttt{Groups} may also have subgroups, which in turn makes the parent \texttt{Group} into a \enquote{metagroup}.
\parencite{ansibleinventoryintro}

The individual \texttt{Hosts} are then declared as children of a specific Ansible \texttt{Group}.
These \texttt{Hosts} need to have at least a name and hostname, so Ansible can access them.
Ansible will default to the declared inventory name as the hostname if no hostname is specified.

The format for specifying an inventory can be found below:

%TC:ignore
\begin{lstlisting}[language=yaml,caption=Example Ansible inventory with \texttt{Groups} and \texttt{Hosts} (source: \cite{ansibleinventoryintro})]
leafs:
  hosts:
    leaf01:
      ansible_host: 192.0.2.100
    leaf02:
      ansible_host: 192.0.2.110

spines:
  hosts:
    spine01:
      ansible_host: 192.0.2.120
    spine02:
      ansible_host: 192.0.2.130

network:
  children:
    leafs:
    spines:

webservers:
  hosts:
    webserver01:
      ansible_host: 192.0.2.140
    webserver02:
      ansible_host: 192.0.2.150

datacenter:
  children:
    network:
    webservers:
\end{lstlisting}
%TC:endignore

The listing shows an inventory file declaring the \texttt{Groups} leafs, spines, network, webservers and datacenter, with the corresponding \texttt{Hosts}.
Notably, spines and leafs are subgroups of network, and network and webservers are subgroups of datacenter.

%TODO: Check if we need to explain vars in detail.

\subsection{Ansible \texttt{Playbooks}}
A \texttt{Playbook} in Ansible is a file that contains several \texttt{Plays}, which in turn contain \texttt{Tasks} to perform on a target.
The \texttt{Tasks} here are operations like copying a file or setting the timezone for a host.
The \texttt{Playbook} below is an example of how one could manage web servers and database servers using an Ansible \texttt{Playbook}:

%TC:ignore
\begin{lstlisting}[language=yaml,caption=Ansible \texttt{Playbook} example (source: \cite{ansibleplaybooks})]
---
- name: Update web servers
  hosts: webservers
  remote_user: root

  tasks:
  - name: Ensure apache is at the latest version
    ansible.builtin.yum:
      name: httpd
      state: latest

  - name: Write the apache config file
    ansible.builtin.template:
      src: /srv/httpd.j2
      dest: /etc/httpd.conf

- name: Update db servers
  hosts: databases
  remote_user: root

  tasks:
  - name: Ensure postgresql is at the latest version
    ansible.builtin.yum:
      name: postgresql
      state: latest

  - name: Ensure that postgresql is started
    ansible.builtin.service:
      name: postgresql
      state: started
\end{lstlisting}
%TC:endignore
This \texttt{Playbook} contains two \texttt{Plays} which target the \texttt{Groups} webservers and databases respectively.
On webservers, the first \texttt{Play} ensures the httpd package is installed with the latest version, and then writes the configuration for it to \enquote{/etc/httpd.conf}, templating it from a local template file.
The second \texttt{Play} targets the databases and ensures PostgreSQL is installed and started.
\parencite{ansibleplaybooks}

\subsection{Ansible \texttt{Roles}}
\texttt{Roles} are reusable units of \texttt{Tasks} which can be used to split configurations into reusable chunks to, for example, apply a similar configuration repeatedly to several logical targets without having to contain them in \texttt{Playbooks}.
\texttt{Roles} also allow for easy delivery of templates and other files together with these \texttt{Tasks}.
Once created, \texttt{Roles} can be referenced either in \texttt{Playbooks} or within other \texttt{Roles} as dependencies.
\texttt{Roles} that are required for the execution of a \texttt{Playbook} can be declared in a file called \enquote{requirements.yml}, which can in turn be used to install all \texttt{Roles} listed in the file.
\parencite{ansibleroles, ansiblegalaxyuserguide}
